\item[\textbf{Problem 4}]
Alice designs a new double-DES scheme. The scheme will first DES-encrypt a
message using $K_{1}$ to get an intermediate ciphertext, then
DES-\textbf{decrypt} the intermediate ciphertext using $K_{2}$ to get the final
ciphertext. Is there any vulnerability in Alice's design?

\Solution

Similarly to Problem 2, Alice's design is highly vulnerable to the
meet-in-the-middle attack.
The ciphertext is:
\[
    C = D_{K_{2}}(E_{K_{1}}(P)).
\]

While the total key space is $112$ bits, the effective security of this scheme
is reduced to roughly $2^{57}$ operations, much like standard double-DES. This
vulnerability exists because the algebraic structure allows an attacker to
isolate the intermediate value $I$. 

By rearranging the encryption equation, we see that $E_{K_{2}}(C) = E_{K_{1}}(P)$.
An attacker can:
\begin{enumerate}
    \item Encrypt the plaintext $P$ forward using all possible values of $K_{1}$
          ($I = E_{K_{1}}(P)$) and store the results in a table.
    
    \item Encrypt the ciphertext $C$ "backward" (since the inverse of Alice's
          second step, $D_{K_{2}}$, is $E_{K_{2}}$) using all possible values of
          $K_{2}$ ($I = E_{K_{2}}(C)$).
\end{enumerate}

Instead of an exhaustive search of $2^{112}$ keys, the attacker only needs to
perform two separate searches of $2^{56}$ operations to find a collision at $I$.
Ultimately, replacing the second encryption with a decryption does not mitigate
the fundamental vulnerability that allows the keys to be attacked independently.
